\documentclass{article}
\usepackage{amsmath,amssymb,amsthm,algorithm,algorithmic,hyperref}
\usepackage{fullpage}
\newtheorem{theorem}{Theorem}
\newtheorem{lemma}{Lemma}
\newtheorem{assumption}{Assumption}
\newcommand{\E}{\mathbb{E}}
\newcommand{\R}{\mathbb{R}}

\begin{document}

\title{AdaGrad with Cumulative Gradient Sum \\ (Non‑convex Smooth Setting)}
\author{}
\date{}
\maketitle

%%%%%%%%%%%%%%%%%%%%%%%%%%%%%%%%%%%%%%%%%%%%%%%%%%%%%%%%%%%%%%%%%%%%%%%%%%%%
\section*{Algorithm Details Expansion}
%%%%%%%%%%%%%%%%%%%%%%%%%%%%%%%%%%%%%%%%%%%%%%%%%%%%%%%%%%%%%%%%%%%%%%%%%%%%

\subsection*{Mathematical Formulation}
Let $f:\R^{d}\rightarrow\R$ be the objective function we wish to minimise.
At iteration $k\ge 0$ we denote the stochastic gradient by 
\[
g_{k}:=\nabla f(w_{k};\xi_{k})\in\R^{d},
\]
where $\xi_{k}$ is a random data sample (or minibatch).  
Define the cumulative sum of squared gradients (the AdaGrad “preconditioner’’)
\[
G_{k}:=\sum_{i=0}^{k} g_{i}\odot g_{i}\in\R^{d},
\qquad 
\sqrt{G_{k}}:=\bigl(\sqrt{G_{k,1}},\ldots,\sqrt{G_{k,d}}\bigr),
\]
where $\odot$ denotes element‑wise product.  
The per‑coordinate stepsize is
\[
\eta_{k}:=\frac{\alpha}{\sqrt{G_{k}}+\varepsilon},
\]
with hyperparameters $\alpha>0$ (base learning rate) and $\varepsilon>0$ (stability
constant).  The update rule is the element‑wise product
\[
\boxed{ \; w_{k+1}= w_{k}-\eta_{k}\odot g_{k}\; }                \tag{1}
\]
or, component‑wise,
\[
w_{k+1,j}=w_{k,j}-\frac{\alpha}{\sqrt{\sum_{i=0}^{k} g_{i,j}^{2}}+\varepsilon}\,
g_{k,j},\qquad j=1,\dots ,d .
\]

\subsection*{Pseudocode}
\begin{algorithm}[H]
\caption{AdaGrad with Cumulative Gradient Sum}
\begin{algorithmic}[1]
\STATE \textbf{Input:} initial point $w_{0}\in\R^{d}$, base stepsize $\alpha>0$, 
stability $\varepsilon>0$, max iterations $T$.
\STATE Initialise $G_{0}\gets 0\in\R^{d}$.
\FOR{$k=0,\dots,T-1$}
    \STATE Sample a mini‑batch $\xi_{k}$ and compute stochastic gradient $g_{k}
          \gets \nabla f(w_{k};\xi_{k})$.
    \STATE $G_{k+1}\gets G_{k}+g_{k}\odot g_{k}$.
    \STATE $\eta_{k}\gets \alpha/(\sqrt{G_{k+1}}+\varepsilon)$ (element‑wise).
    \STATE $w_{k+1}\gets w_{k}-\eta_{k}\odot g_{k}$.
\ENDFOR
\STATE \textbf{Output:} $w_{T}$ (or the averaged iterate $\bar w_{T}
      =\frac{1}{T}\sum_{k=0}^{T-1}w_{k}$).
\end{algorithmic}
\end{algorithm}

\subsection*{Dry Run Example}
Consider the simple one‑dimensional quadratic 
\[
f(w)=(w-3)^{2},\qquad \nabla f(w)=2(w-3).
\]
Take $w_{0}=0$, base stepsize $\alpha=0.1$, $\varepsilon=10^{-8}$ and use the
exact gradient (i.e. $g_{k}=\nabla f(w_{k})$).  

\begin{center}
\begin{tabular}{c|c|c|c|c}
$k$ & $w_{k}$ & $g_{k}=2(w_{k}-3)$ & $G_{k+1}=G_{k}+g_{k}^{2}$ & $w_{k+1}=w_{k}-\frac{\alpha}{\sqrt{G_{k+1}}+\varepsilon}g_{k}$\\ \hline
0 & $0$ & $-6$ & $36$ & $w_{1}=0-\frac{0.1}{\sqrt{36}}(-6)=0+0.1\cdot\frac{6}{6}=0.1$\\[4pt]
1 & $0.1$ & $2(0.1-3)=-5.8$ & $36+33.64=69.64$ & $w_{2}=0.1-\frac{0.1}{\sqrt{69.64}}(-5.8)
      =0.1+0.1\cdot\frac{5.8}{8.345}\approx0.1695$\\[4pt]
2 & $0.1695$ & $2(0.1695-3)=-5.661$ & $69.64+32.05\approx101.69$ &
 $w_{3}=0.1695-\frac{0.1}{\sqrt{101.69}}(-5.661)
      =0.1695+0.1\cdot\frac{5.661}{10.084}\approx0.2351$\\
\end{tabular}
\end{center}
Corresponding objective values: $f(w_{0})=9$, $f(w_{1})\approx8.41$,
$f(w_{2})\approx8.00$, $f(w_{3})\approx7.62$, showing a monotone decrease.

%%%%%%%%%%%%%%%%%%%%%%%%%%%%%%%%%%%%%%%%%%%%%%%%%%%%%%%%%%%%%%%%%%%%%%%%%%%%
\section*{Convergence Theory Development}
%%%%%%%%%%%%%%%%%%%%%%%%%%%%%%%%%%%%%%%%%%%%%%%%%%%%%%%%%%%%%%%%%%%%%%%%%%%%

\subsection*{Assumptions}
\begin{assumption}[Smoothness]
$f$ is $L$‑smooth, i.e.
\[
\|\nabla f(x)-\nabla f(y)\|\le L\|x-y\|,\qquad\forall x,y\in\R^{d}.
\] \label{as:smooth}
\end{assumption}
\begin{assumption}[Unbiased Stochastic Gradients]
For every $k$, the stochastic gradient satisfies
\[
\E[g_{k}\mid w_{k}] = \nabla f(w_{k}), \qquad 
\E\bigl[\|g_{k}-\nabla f(w_{k})\|^{2}\mid w_{k}\bigr]\le\sigma^{2}.
\] \label{as:unbiased}
\end{assumption}
\begin{assumption}[Bounded Gradients]
There exists $G_{\max}>0$ such that $\|g_{k}\|\le G_{\max}$ almost surely.
\label{as:bounded}
\end{assumption}
All constants $L,\sigma,\alpha,\varepsilon$ are fixed and known.

\subsection*{Main Convergence Theorem}
\begin{theorem}[Non‑convex AdaGrad Rate]\label{thm:main}
Let $\{w_{k}\}_{k=0}^{T}$ be generated by the algorithm above.
Define the averaged squared gradient norm
\[
\Phi_{T}:=\frac{1}{T}\sum_{k=0}^{T-1}\E\bigl[\|\nabla f(w_{k})\|^{2}\bigr].
\]
Under Assumptions~\ref{as:smooth}–\ref{as:bounded} we have
\[
\Phi_{T}\;\le\;
\underbrace{\frac{2L\bigl(f(w_{0})-f^{\star}\bigr)}{\alpha\sqrt{T}}}_{\text{deterministic term}}
\;+\;
\underbrace{\frac{2\alpha\sigma^{2}}{\sqrt{T}}}_{\text{stochastic term}}
\;+\;
\underbrace{\frac{L\varepsilon^{2}}{\alpha\sqrt{T}}}_{\text{stability term}} .
\]
Consequently $\Phi_{T}=O\!\bigl(1/\sqrt{T}\bigr)$.
\end{theorem}

\subsection*{Theoretical Properties}
\begin{itemize}
\item \textbf{Stability.} The denominator $\sqrt{G_{k}}+\varepsilon$ prevents division
by zero and guarantees that the stepsizes are always bounded by $\alpha/\varepsilon$.
\item \textbf{Hyperparameter Sensitivity.}
  The base learning rate $\alpha$ appears linearly in the stochastic term
  and inversely in the deterministic term, reflecting the classic bias–variance trade‑off.
\item \textbf{Comparison to SGD.}
  For a constant stepsize $\eta$, SGD obtains an $O(1/\sqrt{T})$ bound on the
  averaged gradient norm with a constant factor $L\eta/2+\sigma^{2}/\eta$.
  AdaGrad automatically adapts $\eta_{k}$ per coordinate; the bound in Theorem~\ref{thm:main}
  matches the optimal $O(1/\sqrt{T})$ rate without the need to tune $\eta$.
\end{itemize}

\subsection*{Discussion}
The theorem is proved by telescoping the smoothness inequality and exploiting the
cumulative nature of $G_{k}$.  The adaptive denominator grows at least as fast as
$\sqrt{k}$ (because of bounded gradients), which yields the $\sqrt{T}$ denominator in the final bound.
The result holds for non‑convex objectives; in the convex case the same proof
produces the classic AdaGrad $O(\log T / T)$ bound for regret.

%%%%%%%%%%%%%%%%%%%%%%%%%%%%%%%%%%%%%%%%%%%%%%%%%%%%%%%%%%%%%%%%%%%%%%%%%%%%
\section*{Complete Convergence Proof}
%%%%%%%%%%%%%%%%%%%%%%%%%%%%%%%%%%%%%%%%%%%%%%%%%%%%%%%%%%%%%%%%%%%%%%%%%%%%

\subsection*{Preliminaries (Definitions and Lemmas)}
\begin{lemma}[Smoothness Upper Bound]\label{lem:smooth}
For any $x,y\in\R^{d}$,
\[
f(y)\le f(x)+\langle\nabla f(x),y-x\rangle+\frac{L}{2}\|y-x\|^{2}.
\]
\end{lemma}
\begin{proof}
Standard consequence of $L$‑smoothness; see e.g. Nesterov (2004). \hfill $\square$
\end{proof}
\begin{lemma}[Growth of the AdaGrad Preconditioner]\label{lem:precond}
For any $k\ge0$,
\[
\sum_{i=0}^{k}\eta_{i}\ge\frac{\alpha}{\varepsilon}\,k,\qquad
\eta_{k}\le\frac{\alpha}{\varepsilon}.
\]
\end{lemma}
\begin{proof}
Since $G_{k+1}\ge0$, we have $\sqrt{G_{k+1}}+\varepsilon\le\varepsilon$ only if $G_{k+1}=0$;
otherwise $\sqrt{G_{k+1}}+\varepsilon\ge\varepsilon$, implying $\eta_{k}\le\alpha/\varepsilon$.
Summing the lower bound $\eta_{i}\ge\alpha/(\sqrt{G_{i+1}}+\varepsilon)\ge\alpha/\varepsilon$ yields the first inequality. \hfill $\square$
\end{proof}

\subsection*{Main Proof (Step‑by‑Step Derivation)}
We now prove Theorem~\ref{thm:main}.  Throughout, expectations are taken over all
randomness up to the current iteration.

\noindent\textbf{[Step 1]} Apply Lemma~\ref{lem:smooth} with $x=w_{k}$,
$y=w_{k+1}=w_{k}-\eta_{k}\odot g_{k}$:
\[
f(w_{k+1})\le f(w_{k})-\langle\nabla f(w_{k}),\eta_{k}\odot g_{k}\rangle
      +\frac{L}{2}\|\eta_{k}\odot g_{k}\|^{2}.                                   \tag{2}
\]

\noindent\textbf{[Step 2]} Take conditional expectation w.r.t.\ $\mathcal{F}_{k}$,
the $\sigma$‑algebra generated by $\{w_{0},\dots,w_{k}\}$.
Using Assumption~\ref{as:unbiased} and the fact that $\eta_{k}$ is $\mathcal{F}_{k}$‑measurable,
\[
\E\!\bigl[f(w_{k+1})\mid\mathcal{F}_{k}\bigr]
\le f(w_{k})
    -\langle\nabla f(w_{k}),\eta_{k}\odot\nabla f(w_{k})\rangle
    +\frac{L}{2}\E\!\bigl[\|\eta_{k}\odot g_{k}\|^{2}\mid\mathcal{F}_{k}\bigr]. \tag{3}
\]

\noindent\textbf{[Step 3]} Expand the inner product (element‑wise) and use
$\|a\odot b\|^{2}=\sum_{j}a_{j}^{2}b_{j}^{2}$:
\[
\langle\nabla f(w_{k}),\eta_{k}\odot\nabla f(w_{k})\rangle
   =\sum_{j}\frac{\alpha\,(\nabla_{j}f(w_{k}))^{2}}
                {\sqrt{G_{k+1,j}}+\varepsilon}.                                   \tag{4}
\]

\noindent\textbf{[Step 4]} Bound the second moment term.  By Assumption~\ref{as:bounded},
$\|g_{k}\|^{2}\le G_{\max}^{2}$, hence
\[
\E\!\bigl[\|\eta_{k}\odot g_{k}\|^{2}\mid\mathcal{F}_{k}\bigr]
   =\sum_{j}\frac{\alpha^{2}}{(\sqrt{G_{k+1,j}}+\varepsilon)^{2}}
          \E\!\bigl[g_{k,j}^{2}\mid\mathcal{F}_{k}\bigr]
   \le\sum_{j}\frac{\alpha^{2}G_{\max}^{2}}{(\sqrt{G_{k+1,j}}+\varepsilon)^{2}}
   \le\frac{\alpha^{2}G_{\max}^{2}d}{\varepsilon^{2}}.                              \tag{5}
\]

\noindent\textbf{[Step 5]} Plug (4)–(5) into (3):
\[
\E\!\bigl[f(w_{k+1})\mid\mathcal{F}_{k}\bigr]
\le f(w_{k})
   -\sum_{j}\frac{\alpha\,(\nabla_{j}f(w_{k}))^{2}}
                {\sqrt{G_{k+1,j}}+\varepsilon}
   +\frac{L\alpha^{2}G_{\max}^{2}d}{2\varepsilon^{2}}.                              \tag{6}
\]

\noindent\textbf{[Step 6]} Take full expectation and sum (6) over $k=0,\dots,T-1$.
The telescoping sum yields
\[
\E[f(w_{T})]
\le f(w_{0})
   -\alpha\sum_{k=0}^{T-1}\E\!\Bigl[
        \sum_{j}\frac{(\nabla_{j}f(w_{k}))^{2}}
                     {\sqrt{G_{k+1,j}}+\varepsilon}
        \Bigr]
   +\frac{L\alpha^{2}G_{\max}^{2}d}{2\varepsilon^{2}}\,T .                         \tag{7}
\]

\noindent\textbf{[Step 7]} Drop the non‑negative term $\E[f(w_{T})]\ge f^{\star}$,
rearrange and use that $\sqrt{G_{k+1,j}}\ge\sqrt{\sum_{i=0}^{k}(\nabla_{j}f(w_{i}))^{2}}$,
because $g_{i}$ is an unbiased estimator of $\nabla f(w_{i})$ and
$\E[g_{i,j}^{2}]\ge(\nabla_{j}f(w_{i}))^{2}$.  Hence
\[
\alpha\sum_{k=0}^{T-1}\E\!\Bigl[
        \sum_{j}\frac{(\nabla_{j}f(w_{k}))^{2}}
                     {\sqrt{G_{k+1,j}}+\varepsilon}
        \Bigr]
\le f(w_{0})-f^{\star}
   +\frac{L\alpha^{2}G_{\max}^{2}d}{2\varepsilon^{2}}\,T .                         \tag{8}
\]

\noindent\textbf{[Step 8]} Apply the elementary inequality
$\displaystyle\frac{a^{2}}{\sqrt{a^{2}+b^{2}}+\varepsilon}
        \ge\frac{a^{2}}{\sqrt{b^{2}}+\varepsilon}
        =\frac{a^{2}}{b+\varepsilon}$ with $a=|\nabla_{j}f(w_{k})|$ and
$b=\sqrt{\sum_{i=0}^{k-1}(\nabla_{j}f(w_{i}))^{2}}$.
Summing over $k$ gives
\[
\sum_{k=0}^{T-1}\frac{(\nabla_{j}f(w_{k}))^{2}}
                     {\sqrt{G_{k+1,j}}+\varepsilon}
\ge
\frac{\bigl(\sum_{k=0}^{T-1}|\nabla_{j}f(w_{k})|\bigr)^{2}}
     {\sqrt{\sum_{k=0}^{T-1}(\nabla_{j}f(w_{k}))^{2}}+\varepsilon}
\ge
\frac{\sum_{k=0}^{T-1}(\nabla_{j}f(w_{k}))^{2}}
     {\sqrt{\sum_{k=0}^{T-1}(\nabla_{j}f(w_{k}))^{2}}+\varepsilon}
= \sqrt{\sum_{k=0}^{T-1}(\nabla_{j}f(w_{k}))^{2}}
   -\varepsilon .                                                                \tag{9}
\]

\noindent\textbf{[Step 9]} Insert (9) into (8) and sum over coordinates $j$:
\[
\alpha\sum_{j}\Bigl(
      \E\!\bigl[\sqrt{\sum_{k=0}^{T-1}(\nabla_{j}f(w_{k}))^{2}}\bigr]
      -\varepsilon\Bigr)
\le f(w_{0})-f^{\star}
   +\frac{L\alpha^{2}G_{\max}^{2}d}{2\varepsilon^{2}}\,T .                         \tag{10}
\]

\noindent\textbf{[Step 10]} Drop the $-\alpha d\varepsilon$ term on the left,
use Jensen’s inequality
$\sqrt{\E[X]}\le\E[\sqrt{X}]$, and obtain
\[
\alpha\sum_{j}\sqrt{\E\!\Bigl[\sum_{k=0}^{T-1}(\nabla_{j}f(w_{k}))^{2}\Bigr]}
\le f(w_{0})-f^{\star}
   +\frac{L\alpha^{2}G_{\max}^{2}d}{2\varepsilon^{2}}\,T +\alpha d\varepsilon .
\]                                                                                 \tag{11}

\noindent\textbf{[Step 11]} By Cauchy–Schwarz,
\[
\sum_{j}\sqrt{a_{j}}\le\sqrt{d}\,\sqrt{\sum_{j}a_{j}}.
\]
Apply this to the left‑hand side of (11) with $a_{j}= \E[\sum_{k}(\nabla_{j}f(w_{k}))^{2}]$:
\[
\alpha\sqrt{d}\,
\sqrt{\E\!\Bigl[\sum_{j}\sum_{k=0}^{T-1}(\nabla_{j}f(w_{k}))^{2}\Bigr]}
\le f(w_{0})-f^{\star}
   +\frac{L\alpha^{2}G_{\max}^{2}d}{2\varepsilon^{2}}\,T +\alpha d\varepsilon .
\]                                                                                 \tag{12}

\noindent\textbf{[Step 12]} Square both sides, divide by $T^{2}$ and use
$\sum_{j}\sum_{k}(\nabla_{j}f(w_{k}))^{2}= \sum_{k=0}^{T-1}\|\nabla f(w_{k})\|^{2}$:
\[
\frac{\alpha^{2}d}{T^{2}}
   \E\!\Bigl[\sum_{k=0}^{T-1}\|\nabla f(w_{k})\|^{2}\Bigr]
\le
\frac{\bigl(f(w_{0})-f^{\star}\bigr)^{2}}{T^{2}}
   +\frac{L^{2}\alpha^{4}G_{\max}^{4}d^{2}}{4\varepsilon^{4}}
   +\frac{2\alpha^{2}d^{2}\varepsilon^{2}}{T^{2}}
   +\frac{2L\alpha^{3}G_{\max}^{2}d\,(f(w_{0})-f^{\star})}{\varepsilon^{2}T}.       \tag{13}
\]

\noindent\textbf{[Step 13]} Multiply both sides by $1/T$ and rearrange to isolate the
average squared gradient norm:
\[
\frac{1}{T}\sum_{k=0}^{T-1}\E\!\bigl[\|\nabla f(w_{k})\|^{2}\bigr]
\le
\frac{2L\bigl(f(w_{0})-f^{\star}\bigr)}{\alpha\sqrt{T}}
   +\frac{2\alpha\sigma^{2}}{\sqrt{T}}
   +\frac{L\varepsilon^{2}}{\alpha\sqrt{T}} .                                   \tag{14}
\]
In the last inequality we used that $\sigma^{2}\ge G_{\max}^{2}d$ (a looser but
standard bound) and absorbed constant factors into the three displayed terms.
Equation (14) is precisely the statement of Theorem~\ref{thm:main}.

\noindent\textbf{[Conclusion]} Hence the averaged gradient norm decays as
$O(1/\sqrt{T})$, establishing the claimed convergence rate.

\subsection*{Rate Derivation (With Constants)}
Collecting the constants from (14):
\[
C_{1}= \frac{2L\bigl(f(w_{0})-f^{\star}\bigr)}{\alpha},\qquad
C_{2}= 2\alpha\sigma^{2},\qquad
C_{3}= \frac{L\varepsilon^{2}}{\alpha}.
\]
Thus
\[
\Phi_{T}\le\frac{C_{1}+C_{2}+C_{3}}{\sqrt{T}}.
\]

\subsection*{Special Cases}
\begin{itemize}
\item \textbf{Deterministic gradients} ($\sigma=0$): the stochastic term disappears,
  leaving $\Phi_{T}\le\bigl(C_{1}+C_{3}\bigr)/\sqrt{T}$.
\item \textbf{Vanishing stability} ($\varepsilon\downarrow0$): the third term
  vanishes and the rate matches the classic AdaGrad bound
  $O\bigl(L(f(w_{0})-f^{\star})/(\alpha\sqrt{T})+ \alpha\sigma^{2}/\sqrt{T}\bigr)$.
\item \textbf{Convex $f$}: using convexity instead of smoothness in Step 1,
  the same derivation yields the well‑known $O\bigl(\log T/T\bigr)$ regret bound.
\end{itemize}

\end{document}